\chapter{Infrastructure as Code: Terraform and Database Change Management}
\index{CI/CD}\index{Terraform}\index{infrastructure as code}\index{schema migration}\index{expand-contract}\index{GitOps}\index{Atlas Terraform provider}\index{database change management}%

\begin{objectives}
\begin{itemize}
    \item Provision Atlas projects, clusters, users, and private endpoints with the MongoDB Atlas Terraform provider.
    \item Manage search/vector indexes, roles, and network resources as reviewed code.
    \item Automate index and \texttt{\$jsonSchema} migrations as numbered, idempotent CI/CD artifacts.
    \item Execute zero-downtime expand--migrate--contract schema evolution with feature-flagged read switches.
    \item Build the multi-environment GitOps pipeline in the Thursday lab.
    \item Architect change management where infrastructure, indexes, and schemas ship as one reviewed unit.
\end{itemize}
\end{objectives}

\begin{crosscloud}
Infrastructure as code is the most vendor-neutral skill in this book: Terraform manages all six platforms, and the GitOps discipline---plan in CI, apply on merge, state locked remotely---is identical everywhere.

\vspace{2mm}
{\footnotesize
\begin{tabular}{@{}p{2.9cm}p{3.0cm}p{3.0cm}p{3.0cm}@{}} 
\toprule
\textbf{Capability} & \textbf{AWS} & \textbf{Azure} & \textbf{GCP} \\
\midrule
Native IaC & CloudFormation / CDK & ARM / Bicep & Deployment Manager \\
Terraform provider & \texttt{hashicorp/aws} & \texttt{azurerm} & \texttt{google} \\
DB migration tools & Flyway / Liquibase / Alembic & Same & Same \\
Plan-in-CI gate & \texttt{terraform plan} on PR & Same & Same \\
State backend & S3 + DynamoDB lock & Azure Blob + lease & GCS + lock \\
Secrets & Secrets Manager + CI vars & Key Vault + CI vars & Secret Manager + CI vars \\
\addlinespace
Snowflake / Fabric & \multicolumn{3}{l}{Snowflake provider for accounts/warehouses/grants; schemachange for migrations; Fabric via Terraform azapi and Git integration} \\
\bottomrule
\end{tabular}}

\vspace{1mm}
MongoDB's distinctive gap and opportunity: schema is not enforced by the engine, so \emph{change management}---validators, indexes, migrations---is the mechanism that keeps forty services' implicit schemas from drifting apart. CI/CD is where that discipline lives \cite{mongodbTerraformDocs,terraformDocs}.
\end{crosscloud}

\section{Week 21 Roadmap: The Database in the Pipeline}
Everything prior in this book was learned against systems you built by hand; production forbids that. The week's claim: any change to infrastructure, indexes, roles, or validation rules that is not a reviewed, versioned, automatically applied artifact will eventually be an incident. We build the pipeline that makes the reviewed artifact the path of least resistance \cite{mongodbTerraformDocs}.

\begin{definitionbox}
\textbf{GitOps for databases} is the operating model where the Git repository is the desired state of infrastructure and schema, CI validates proposed changes (\texttt{terraform plan}, migration dry-runs, contract tests), and CD applies merged changes to environments in order (dev $\to$ staging $\to$ prod) with gates.
\end{definitionbox}

\begin{keyterms}
\textbf{Terraform state}: the recorded mapping between configuration and real resources; stored remotely with locking. \textbf{Migration}: a numbered, idempotent script changing schema or indexes. \textbf{Expand--contract}: the zero-downtime schema evolution pattern (add new, dual-write, backfill, switch reads, remove old). \textbf{Contract test}: the CI check asserting a service's schema assumptions still hold. \textbf{Drift}: reality diverging from the repository.
\end{keyterms}

\section{Monday: The Atlas Terraform Provider}
\subsection{What the Provider Owns}
The \texttt{mongodbatlas} provider manages the Atlas control plane: organizations, projects, clusters (including sharded topologies and region configs), database users and custom roles, network peering and private endpoints, alert configurations, backup policies, Online Archive rules, and search/vector index definitions. The design consequence: everything from Parts III--V that was ``configured in the UI'' becomes reviewable HCL---zones, alerts, archive rules, QE-adjacent key policies---with the same pull-request rigor as application code \cite{mongodbTerraformDocs}.

\begin{lstlisting}[language=bash, caption={Atlas cluster as Terraform HCL (from the syllabus, annotated).}]
resource "mongodbatlas_cluster" "analytics" {
  project_id   = var.atlas_project_id
  name         = "analytics-sharded"
  cluster_type = "SHARDED"

  replication_specs {
    num_shards = 2
    regions_config {
      region_name     = "US_EAST_1"
      electable_nodes = 3
      priority        = 7
      read_only_nodes = 0
    }
  }

  provider_name               = "AWS"
  provider_instance_size_name = "M30"
  mongo_db_major_version      = "8.0"
  backup_enabled              = true        # Chapter 20/23 control, in code
  pit_enabled                 = true        # PITR window, in code

  labels {
    key   = "environment"
    value = "production"
  }
}
\end{lstlisting}

\subsection{State, Secrets, and the Plan Gate}
Three rules make Terraform safe: state lives in a remote backend with locking (two concurrent applies corrupt state; locking prevents the second), secrets never appear in code (Atlas API keys come from CI secret stores into environment variables), and \texttt{terraform plan} runs on every pull request with the plan output posted for review---the reviewer approves the plan, not just the diff. Drift detection is a scheduled \texttt{plan} expecting zero changes; non-zero drift pages the platform team, because drift means someone changed production outside the repository.

\begin{pitfall}
\textbf{Clicking in the Atlas UI ``just this once.''} Every console change is drift the next apply will either revert (surprising the clicker) or absorb silently (losing the review trail). Emergency console changes are followed same-day by a PR codifying them, or by a deliberate revert---recorded either way.
\end{pitfall}

\section{Tuesday: Indexes, Roles, and Search as Code}
The provider plus scripted \texttt{mongosh} migrations cover the database plane: search and vector index definitions as Terraform resources (Chapters 13--14's indexes reviewed like any schema), custom roles and users as resources (Chapter 17's grant matrix enforced), and alert policies as resources (Chapter 23's alerts-as-code). What Terraform does not own---collection-local artifacts like validators and non-search indexes---ships as numbered migrations (Wednesday). The boundary rule: if the Atlas API owns it, Terraform owns it; if the mongod owns it, migrations own it; nothing is owned by a human clicking.

\section{Wednesday: Automated Migrations}
\subsection{Numbered, Idempotent, Forward-Only}
A migration is a numbered script (\texttt{0007\_expand\_contract.js}) applied exactly once per environment by a runner that records applied versions in a control collection. Migrations are idempotent (safe to re-run), forward-only (rollbacks are new migrations), and throttled when they touch data (batched, resumable writes---a one-shot full-collection update saturates the primary, as the syllabus failure story warns).

\begin{lstlisting}[language=JavaScript, caption={Idempotent expand-phase migration, run by CI.}]
// migrations/0007_expand_order_total.js
db.orders.updateMany(
  { order_total: { $exists: false }, total_price: { $exists: true } },
  [ { $set: { order_total: "$total_price" } } ] );     // batched + throttled in prod
db.orders.createIndex({ order_total: 1, status: 1 });
db.migration_log.insertOne({ id: "0007", at: new Date(), by: "ci" });
\end{lstlisting}

\subsection{Expand--Migrate--Contract in Full}
The zero-downtime rename from the syllabus failure story, done right: \textbf{expand}---deploy application code writing both \texttt{total\_price} and \texttt{order\_total}; \textbf{migrate}---backfill existing documents with a throttled, resumable pipeline; \textbf{switch reads}---flip consumers to the new field behind a feature flag, with a compatibility view shim for lagging consumers; \textbf{contract}---only when zero lagging readers remain, enforce the new field with \texttt{\$jsonSchema} (\texttt{validationAction} graduating \texttt{warn} $\to$ \texttt{error}) and \texttt{\$unset} the old field in a final batched migration. Every step is a separate reviewed migration; the failure story happened precisely because backfill and read-switch shipped in one release.

\begin{notebox}
The Week 3 schema contract test belongs in this same pipeline: a CI job asserting that a service's expected document shape still validates against the collection's \texttt{\$jsonSchema}, so schema, index, and infrastructure changes deploy as one reviewed unit.
\end{notebox}

\section{Thursday Hands-On Project: The Three-Environment GitOps Pipeline}
Build the pipeline: Terraform modules for dev/staging/prod Atlas environments, a migrations runner, and CI that plans on PR and applies on merge per environment.

\subsection{Repository Layout and Pipeline}
\begin{lstlisting}[language=bash, caption={GitOps repository structure and CI stage outline.}]
repo/
  infra/            # Terraform: modules + envs (dev, staging, prod)
    envs/prod/main.tf
    modules/cluster/
  migrations/       # numbered mongosh scripts + runner
    0006_orders_indexes.js
    0007_expand_order_total.js
  k8s/              # Chapter 22 operator manifests (staging)
  .github/workflows/db-change.yml

# CI stages:
#   PR open      -> terraform plan (all envs) + migration lint + contract tests
#   merge        -> apply dev -> smoke tests -> apply staging -> smoke tests
#   manual gate  -> apply prod -> post-deploy verification -> tag release
\end{lstlisting}

The lab exercises the failure modes deliberately: a plan showing an unexpected destroy is caught at PR review; a staging smoke-test failure blocks prod; and the expand--contract sequence for \texttt{total\_price} $\to$ \texttt{order\_total} runs end-to-end with the old app version still reading during the rollout---proving no orders are dropped, in contrast to the syllabus's cautionary tale.

\subsection*{Deliverables}
\begin{itemize}
    \item The repository with Terraform modules, two migrations, and the CI workflow file.
    \item Evidence of the plan gate catching a deliberately bad change.
    \item The expand--contract run with dual-version application traffic and zero dropped orders.
\end{itemize}

\subsection*{Acceptance Criteria}
\begin{itemize}
    \item A PR produces plan output for all environments without applying anything.
    \item Staging failure demonstrably blocks the production apply.
    \item The migration runner applies each migration exactly once and is idempotent on re-run.
\end{itemize}

\subsection*{Stretch Goals}
\begin{itemize}
    \item Add scheduled drift detection (zero-change \texttt{plan}) with alerting on drift.
    \item Implement the contract phase as a time-delayed migration gated on a ``no lagging readers'' check.
\end{itemize}

\section{Friday Use Case and Architecture: Change Management for a Regulated Platform}
\subsection{Scenario and Requirements}
A fintech with SOX-adjacent change-control obligations must prove that every production database change was reviewed, tested in staging, applied by automation (not a human), and is attributable to an approved change request.

\subsection{Architecture}
The pipeline is the control: Git branch protection requires review on \texttt{infra/} and \texttt{migrations/}; CI links each PR to a change ticket; applies run as a CI service account holding least-privilege Atlas API keys (Chapter 17) from the CI vault; staging applies and smoke tests gate production; and the audit pack---PR, plan output, staging evidence, apply log, approver identity---is assembled per change automatically. Atlas audit events (Chapter 17) close the loop: any change to the cluster \emph{not} correlated with a pipeline run within an hour triggers the drift alert and a change-control incident.

\begin{pitfall}
\textbf{A pipeline humans can bypass.} If engineers retain direct Atlas write access ``for emergencies,'' the control is fiction. Human write access is break-glass only (alarmed, vaulted, expiring), and every break-glass use generates a retroactive PR plus a review note.
\end{pitfall}

\section{Interview Q\&A}
\subsection{DataOps}
\begin{enumerate}
    \item \textbf{Q: Why remote state with locking?}\\
    A: State is the source of truth for what exists; concurrent applies corrupt it, and local state is unshareable and unauditable. Locking serializes applies; the remote backend makes them reviewable.

    \item \textbf{Q: Why forward-only migrations?}\\
    A: Down migrations are code paths tested rarely and needed under maximum stress. Rollback = a new forward migration, reviewed and tested like any change.

    \item \textbf{Q: What does expand--contract actually buy?}\\
    A: Every intermediate state is valid for old and new code simultaneously, so rolling deploys never straddle an incompatible schema---the failure mode of same-release rename + backfill.

    \item \textbf{Q: Where do Atlas API keys live?}\\
    A: In the CI secret store, injected at runtime, scoped least-privilege, rotated on schedule---never in code, never in Terraform variable files committed to the repo.

    \item \textbf{Q: Why bundle the schema contract test into the infra pipeline?}\\
    A: Because schema, indexes, and infrastructure are one logical change; separating their validation lets a compatible-looking pair ship an incompatible system.
\end{enumerate}

\section{Further Reading}
The Atlas Terraform provider registry documents every resource used here \cite{mongodbTerraformDocs}, and Terraform's own documentation covers backends and locking \cite{terraformDocs}. JSON Schema validation underpins the contract phase \cite{mongodbJsonSchemaDocs}. The GitOps operating model and its audit properties align with the SRE release-engineering chapters \cite{sreBook}, and Chapter 22 extends the same manifests-and-operators discipline onto Kubernetes \cite{mongodbK8sOperatorDocs}.

\section*{Key Takeaways}
\begin{itemize}
    \item Terraform owns the control plane; numbered migrations own the data plane; humans own neither directly.
    \item Plan-on-PR is the review gate; zero-drift scheduled plans detect out-of-band change.
    \item Migrations are numbered, idempotent, forward-only, and throttled when they touch data.
    \item Expand--migrate--contract makes every intermediate state valid for old and new code.
    \item Secrets come from the CI vault; API keys are least-privilege service accounts.
    \item Auditability is the pipeline's byproduct: PR + plan + staging evidence + apply log per change.
\end{itemize}

\section{Exercises}
\begin{enumerate}
    \item \textbf{(Conceptual)} Draw the boundary line between Terraform-owned and migration-owned resources for your current project.
    \item \textbf{(Conceptual)} Why does the contract phase require zero lagging readers, and how do you prove that count is zero?
    \item \textbf{(Conceptual)} Design the rollback for a destructive contract migration discovered wrong after release.
    \item \textbf{(Hands-on)} Implement the migration runner with the control collection and demonstrate idempotent re-runs.
    \item \textbf{(Hands-on)} Inject drift manually in the Atlas UI and show the scheduled plan detecting it.
    \item \textbf{(Design)} Write the change-control policy satisfying the fintech's audit: gates, evidence, break-glass, retention.
    \item \textbf{(Design)} Design migration throttling: batch size, sleep, oplog-window guard, and the abort condition.
    \item \textbf{(Interview)} ``A junior applied Terraform to prod before staging passed.'' Enumerate the missing gates and the immediate remediation.
\end{enumerate}

\section{Capstone Project: The One-Pipeline Change Platform}\label{sec:ch21-capstone}

\begin{capstonebox}
\textbf{Mission:} deliver the change-management platform end-to-end: Terraform-managed three-environment Atlas infrastructure, a migration runner, contract tests, drift detection, and a demonstrated expand--contract evolution---with the per-change audit pack assembled automatically.

\textbf{Estimated time:} 6--8 hours.

\textbf{What you will practice:}
\begin{itemize}
    \item Structuring multi-environment Terraform with modules and remote state.
    \item Running schema evolution with zero downtime and dual-version traffic.
    \item Producing change-control evidence as a pipeline artifact.
\end{itemize}
\end{capstonebox}

\subsection*{Scenario}
The fintech adopts your pipeline as the mandatory path for all database change. The deliverable must survive an auditor reading one change end-to-end and an engineer executing one emergency change correctly.

\subsection*{Requirements}
\begin{itemize}
    \item Terraform modules + three environments with remote, locked state; no secrets in the repo.
    \item The migration runner (numbered, idempotent, logged) plus at least three real migrations including one expand--contract pair.
    \item CI: plan on PR, contract tests, staged applies with smoke-test gates, manual prod gate.
    \item Scheduled drift detection with a demonstrated detection.
    \item The per-change audit pack assembled by CI and the break-glass procedure documented and drilled.
\end{itemize}

\subsection*{Implementation Steps}
\begin{enumerate}
    \item \textbf{Foundation.} Remote state, modules, environments; least-privilege CI service account.
    \item \textbf{Migrations.} Runner + the three migrations; idempotency and throttle evidence.
    \item \textbf{Pipeline.} Full CI flow with gates; demonstrate a staging failure blocking prod.
    \item \textbf{Evolution.} Execute expand--contract under dual-version traffic; prove zero dropped orders.
    \item \textbf{Audit.} Assemble the audit pack; drill the break-glass path with its retroactive PR.
\end{enumerate}

\subsection*{Deliverables}
\begin{itemize}
    \item The repository, CI configuration, and all evidence artifacts.
    \item The change-control policy document and the drilled break-glass record.
    \item The expand--contract evidence pack.
\end{itemize}

\subsection*{Acceptance Criteria}
\begin{itemize}
    \item A full change flows PR $\to$ plan $\to$ dev $\to$ staging $\to$ gate $\to$ prod with no manual cluster access.
    \item Drift and break-glass use are both detected and generate tracked follow-ups.
    \item The auditor's end-to-end read of one change requires no oral explanation.
\end{itemize}

\subsection*{Stretch Goals}
\begin{itemize}
    \item Add automated rollback: a revert PR generated from the failed change's plan.
    \item Extend contract tests to consumer-driven contracts for the top five services.
\end{itemize}
